% SPDX-License-Identifier: CC-BY-4.0
% Copyright (C) 2026 Fabian Januszewski, Abhijit (the Langlandscape team)
% GENERATED by tools/extract_statements.py from mathematics/chapters/ -- do not edit by hand.

\begin{quoted}{Hypothesis}{18.15}{139}{hyp:s05-X}{$\HypX(M;q)$; \texttt{if:s05-hypX}; seed Hyp.~8.1}{chapters/s05/hypX.tex}
\emph{Data.} A prime $p\ge5$, a prime $q\ne p$, an integer $M\ge1$ with $(M,pq)=1$; $B=B_q$; for $s\ge0$:
$\Gamma_s=\Gamma_1(Mp^s)\cap\Gamma_0(q)$, $X_s$ its compactified modular curve over $\Q$ in the $\mu_N$-model, $J_s=\Jac X_s$ with Picard
functoriality (decided, Fabian, 2026-11-15: FD1), $\Xq(J_s)$ the character group of the toric part of the special fibre at $q$ of its Néron model,
$S_s=\XBq{\Uone{Mp^s}}$, $\Z[S_s]^0$ the degree-zero part of the homological Brandt module, $\cH_s$ and
$\theta_s:\cH_s\to\End\Z[S_s]^0$ as in Def.~18.4 ($\langle d\rangle\mapsto\langle d\rangle^{-1}$, $U_p\mapsto N_{U_p}$ for $s\ge1$,
$T_p\mapsto T_p$ for $s=0$, $U_q\mapsto+\Pi_q$).

\emph{Statement.} There is a family of isomorphisms of abelian groups $\varphi_s:\Xq(J_s)\xrightarrow{\sim}\Z[S_s]^0$, $s\ge0$, such that:
\provisional{S05-M3: one family $(\varphi_s)_s$, as (X3) links $\varphi_s$ and $\varphi_{s'}$; the stage tex's Lean track writes
$\forall s\,\exists\varphi$ (Remark~18.16).}
\begin{enumerate}[label=(X\arabic*)]
\item \emph{(Hecke)} $\varphi_s\circ h=\theta_s(h)\circ\varphi_s$ for all $s\ge0$ and all $h\in\cH_s$;
\item \emph{(pairing)} $\langle\varphi_sx,\varphi_sy\rangle_w=u_s(x,y)$ for all $s\ge0$ and $x,y\in\Xq(J_s)$, where $u_s$ is Grothendieck's
  monodromy pairing, normalised to be positive definite (\pathtt{rem:s05-illusie}), and $\langle[e],[e']\rangle_w=w_e\delta_{ee'}$ with
  $w_e=|\Gamma_e/(\Gamma_e\cap\{\pm1\})|$ (\pathtt{if:s03-graph});
\item \emph{(naturality)} for $s'\ge s\ge0$, let $\beta_{s's}:\Xq(J_{s'})\to\Xq(J_s)$ be the map on character groups induced by the
  Picard pullback $J_s\to J_{s'}$ along the natural projection $X_{s'}\to X_s$ (contravariance of $\Xq$). Then
  $\varphi_s\circ\beta_{s's}=\pi_*\circ\varphi_{s'}$, with $\pi_*:\Z[S_{s'}]^0\to\Z[S_s]^0$ the pushforward along $\pi:S_{s'}\to S_s$. (Direction: the torus of the special fibre of a semistable Jacobian is
  $H^1(\Gamma,\Z)\otimes\mathbb G_m$ for the dual graph $\Gamma$ \lit{mazur\_1977\_modular-curves-and-the-eisenstein-ideal}{Appendix, \S3,
  p.~183, for $N=p$}, so $\Xq=H_1(\Gamma,\Z)$, and the map on characters induced by Picard pullback is dual to pullback on $H^1$, i.e.\
  the pushforward on $H_1$ of the dual graphs; this is why (X3) pairs $\beta_{s's}$ with $\pi_*$. \stsketch{} at $\Gamma_1$-level.)
\item \emph{(Galois)} for all $s\ge0$: on $T_pJ_s$, with $G_{\Q_q}$ and $\cH_s$ acting, the filtration $T_t\subset T^{I_q}\subset T_pJ_s$
  (Thm.~18.30(v), stated for any semistable abelian variety) has, as $\cH_s[G_{\Q_q}]$-modules, the graded pieces
  \[\begin{gathered}\text{bottom }T_t:\ \Hom(\Z[S_s]^0,\Zp(1))\ (\text{transpose action}),\\
  \text{middle }T^{I_q}/T_t:\ T_pJ_1(Mp^s)^{\oplus2}\ (\text{via the two degeneracies; for the operators away from $q$}),\\
  \text{top }T_pJ_s/T^{I_q}:\ \Z[S_s]^0\otimes\Zp\ (\text{action }h\mapsto\theta_s(h^\Ros)),\end{gathered}\]
  where $h^\Ros$ is the Rosati dual (for $T_\ell$: $\langle\ell\rangle^{-1}T_\ell$); for $h$ with $h^\Ros\notin\cH_s$ ($U_p$, and $U_\ell$ for
  $\ell\mid M$, whose Rosati duals involve $w_{Mp^s}$), $\theta_s(h^\Ros)$ \emph{means} the $\langle,\rangle_w$-adjoint $\theta_s(h)^*$ (the two readings
  agree where both are defined, Lemma~18.12(b)); the arithmetic $\Frob_q$ acts on the bottom piece as $q\,U_q$
  (equivalently $\sigma_q=U_q^{-1}$ on $\Xq(J_s)$, i.e.\ the Frobenius permutation of $S_s$ is right multiplication by $\Pi_q^{-1}$ under
  $\varphi_s$; this uses the $\mu_N$-model)
  (Frobenius clause in the $\mu_N$-model with Picard functoriality, as decided (Fabian, 2026-11-15, FD1); with the naive model it would read
  $q\langle q\rangle^{-1}U_q$, Remark~18.10);
  and the monodromy operator $N$ from the top to the bottom piece is $x\mapsto\langle x,\cdot\rangle_w$ (via (X2)). Transported to
  $H^1_{\rm par}(\Gamma_s,\Zp)(1)$ by the comparison of $T_pJ_s$ with Betti cohomology.
\end{enumerate}
$\HypXchar(M;q):=$ (X1)$\wedge$(X2)$\wedge$(X3); $\HypXgal(M;q):=$ (X4) (presupposing the family $\varphi$); $\HypX:=$ both. \sthyp.
At $s=0$ and $M\le2$, the $u$-adjoint of $\beta_{s'0}$ is compared with the transfer $\pi^!$ with the factor $2$ of Remark~18.3 (Lemma~18.12(c)).
\end{quoted}

\begin{quoted}{Proposition}{18.6}{128}{prop:s05-sign}{$\varepsilon=+1$; exponent $+1$}{chapters/s05/sign.tex}
\begin{enumerate}[label=(\alph*)]
\item ($\Gamma_0$-level, weight $2$.) Let $N=qM'$ with $q\nmid M'$, and identify $\Xq(J_0(N))\cong\Z[\Sigma(M')]^0$ (Ribet's
  Prop.~3.1; Thm.~18.30(ii)), and $\Sigma(M')\cong\XBq{\Uzero{M'}}$ by Deuring's correspondence as in Ribet \S3. Then $U_q=T_q$ acts
  on $\Xq(J_0(N))$ as the permutation $[x]\mapsto[x\Pi_q]$; so $\varepsilon=+1$. Since $\Pi_q^2$ acts trivially on
  $\XBq{\Uzero{M'}}$, this permutation equals $[x]\mapsto[x\Pi_q^{-1}]$ (the Frobenius permutation). \stpmc.
  At $\Gamma_1$-level, in the $\mu_N$-model of the Setting, the Frobenius permutation is $[x]\mapsto[x\Pi_q^{-1}]$ and differs from
  $[x]\mapsto[x\Pi_q]$ in general: \stpmc{} given (X1) and FD1 (it follows from (X1) and the four-row table, Remark~18.10, anchored at the naive/Albanese row, which is \stpmc{} conditional on FD1; Ribet--Stein's Lemma~3.17 agrees under a convention).
  (The model is fixed by the decision FD1 (Fabian, 2026-11-15): the Frobenius permutation at $\Gamma_1$-level ($\mu_N$-model, Picard functoriality) is $\Pi_q^{-1}$; the status \stpmc{} above is that of the derivation (given FD1), not of the model choice.)
\item (Consistency with stage 01.) For $D=q=11$, $M'=1$: $\Pi_{11}=\mathrm{id}$ on $\Z[\XB{\hat\cO^\times}]$, and $a_{11}(11a)=+1$
  (Example~5.6); so $\varepsilon=+1$. \stproved{} that $\Pi_{11}=\mathrm{id}$; $\varepsilon=+1$ then follows given the \stknown{} value $a_{11}=+1$ and the compatibility of
  $\theta$ with $U_{11}$ on the cuspidal line ((X1) at $\Gamma_0(11)$, i.e.\ (a), or JL at $11$; Example~5.6 is itself \stpmc).
\item (Exponent.) Let $N\ge1$ with $(N,q)=1$ and let $\cH^{\rm new}$ be the image of the weight-$2$ Hecke algebra of
  $\Gamma_1(N)\cap\Gamma_0(q)$ (generated by the $T_\ell$, $\ell\nmid Nq$, the $U_\ell$, $\ell\mid N$, the classical
  $\langle d\rangle$, $d\in(\Z/N)^\times$, and $U_q$) in the endomorphisms of the $q$-new subspace. Let
  $\theta:\cH^{\rm new}\to\End\Z[\XBq{\Uone N}]^0$ be a ring homomorphism with $\theta(\langle d\rangle)=\langle d\rangle^{-\delta}$ (adelic)
  and $\theta(U_q)=\varepsilon\Pi_q^{\,e}$, $\varepsilon,\delta,e\in\{\pm1\}$. If $\langle q\rangle^2\ne1$ on $\Z[\XBq{\Uone N}]^0\otimes\C$ (equivalently:
  some character $\chi$ of $(\Z/N)^\times$ with $\chi(q)\ne\pm1$ occurs in the adelic diamond action on it), then $e=\delta$. If
  $\langle q\rangle^2=1$ there (e.g.\ $N\le2$, or $q^2\equiv\pm1\bmod N$: only even $\chi$ occur, since $\langle-1\rangle=1$ on the class set,
  see the proof), the relation imposes nothing. \stpmc{}; checked numerically in T05-2, Remark~18.7(2).
  \provisional{S05-M1: compared with the stage tex, (c) constrains only $\theta(\langle d\rangle)$ and $\theta(U_q)$ (the other generators are
  arbitrary), and the condition on $\chi$ is read on the definite side; both are harmless generalisations (self-audit F8).}
\item (The diamond convention is $\delta=+1$, among the dictionaries considered here.) Let $\theta'=(\theta'_s)_{s\ge1}$ be either $(\theta_s)_s$ or $(\theta^\Ros_s)_s$
  (Prop.~18.18). If (X1)--(X3) of Hypothesis~18.15 hold for some family $(\varphi_s)$ with $\theta'$ in place of $\theta$,
  and $\rk_{\Zp}e\,\Zp[S_s]\ge2$ for some $s\ge1$, then $\theta'=\theta$; i.e.\ $\delta=+1$ among $\theta_s$, $\theta^\Ros_s$ and the dictionaries with $T_\ell\mapsto T_\ell$ \emph{as soon as some $\langle\ell\rangle^2\ne1$ on $\Z[S_s]^0$} (at $p=5$, $s=1$, $M=1$ only even $\chi$ with $\chi^2=1$ occur, so this qualifier is needed), hence $e=+1$ whenever (c) applies:
  $U_q\mapsto+\Pi_q$. The dictionaries with $\delta=-1$ and $T_\ell\mapsto T_\ell$ are excluded by the $T_\ell$-relation as soon as some
  $\langle\ell\rangle^2\ne1$ on $\Z[S_s]^0$ (Lemma~18.12(a)). This agrees with stage 01's diamond flag on forms (adelic
  $\langle d\rangle$ $=$ classical $\langle d\rangle^{-1}$, Remark~4.5). \stpmc{} (the input ``Picard $U_p$ commutes with the
  pullbacks'' is Lemma~18.9, \stpmc; batch-5 closure, previously \sttoprove). The hybrid dictionary $\langle d\rangle\mapsto\langle d\rangle$, $T_\ell\mapsto\langle\ell\rangle T_\ell=T_\ell^*$, $U_q\mapsto\Pi_q^{-1}$, $U_p\mapsto N_{U_p}$ ($\delta=-1$) satisfies the
  relations of Lemma~18.12(a), (b) and commutes with $\pi_*$, and is \emph{not} excluded by these arguments (nor by (X3)). Nor is it excluded by T05-2: at the levels computed there ($s=0$, no $U_p$) it
  coincides with $\theta^\Ros_0$, which matches the classical eigensystems as well as $\theta_0$ does (Remark~18.7(2); this is
  Prop.~18.18 at a single level). Excluding it \emph{modulo cited results}: under the hybrid dictionary the Brandt eigenvector attached to an eigenform $f$ carries the eigensystem of $\bar f$ away from $Mpq$, but $U_p\mapsto a_p(f)$ in place of $a_p(\bar f)=\overline{a_p(f)}$. The eigensystem matching away from $Mpq$ is Jacquet--Langlands
  (\lit{jacquet-langlands\_1970\_automorphic-forms-on-gl2}{Ch.~III, \S14, Thm.~14.4, p.~470 (printed pagination; PDF p.~479)}, with strong multiplicity one), which says nothing about $U_p$ at $p\mid$ level. The $U_p$-step needs \emph{local newform theory at $p$}
  (Casselman 1973, \emph{On some results of Atkin and Lehner}, Math.\ Ann.\ 201; not in \pathtt{literature/}, not checked): if $f$ is \emph{new at $p^s$}, the $U_1(p^s)$-fixed space of the local component $\pi_p$ is a line on which $U_p$ acts by $a_p(f)$, and $\bar f$ is again new at $p^s$ with $a_p(\bar f)=\overline{a_p(f)}$; the Brandt side has the same $\GL_2(\Q_p)$-module, so the $U_p$-eigenvalue on the line of $\bar f$ must be $\overline{a_p(f)}$. Hence the hybrid is excluded at every level $s\ge1$ carrying a form $f$ that is $q$-new, \emph{new at $p^s$} and has $a_p(f)\notin\R$ (this forces the $p$-part of the nebentypus to be primitive of conductor $p^s$: for trivial $p$-part $a_p$ is real). The hypothesis ``new at $p$'' cannot be dropped: for a $p$-stabilisation $f$ of a $p$-old form with $\chi(p)=1$, the $U_p$-eigenvalues on the $\bar f$-old space are $\{\bar\alpha,\bar\beta\}=\{p/\alpha,\bar\chi(p)\alpha\}\ni\alpha=a_p(f)$, so no contradiction arises, although $\alpha\notin\R$. Thm.~6.3 for the $\Lambda$-adic form is not re-checked here. \stpmc{} (cited: JL, strong multiplicity one, local newform theory; the existence of such an $f$ at some level is an extra condition we do not verify in general). All of (d) is a statement about the shape of
  $\HypX$, which itself stays \sthyp.
\end{enumerate}
Part (e) of the stage tex (higher weight, local Jacquet--Langlands for $\St\otimes\chi$) is written in the section on higher weight
(WP S05-M2).
\end{quoted}

\begin{quoted}{Proposition}{18.24}{150}{prop:s05-Xiso}{$\HypXiso$ holds non-canonically}{chapters/s05/subquot.tex}
Under the hypotheses of Def.~18.23, assume:
\begin{enumerate}[label=(\roman*)]
\item $H:=eH^1_{\rm par}(\Gamma_s,\Zp)^{\pm}_\m$ is free of rank one over $\T_\m$, and $\T_\m$ is $\Zp$-Gorenstein;
\item the hypotheses of Prop.~18.25 hold for $A=\Zp$, $\T_\m$, $\fa$ and $\fb:=\ker(\T_\m\to\T''_\m)$ ($\T''$ the $q$-old
  quotient): $\T_\m$ is finite free over $\Zp$ (automatic: finite and torsion-free over the PID $\Zp$), $\T'_\m$, $\T''_\m$ are torsion-free
  (automatic: images in endomorphism rings of lattices), $\fa\cap\fb=0$ and $(\T_\m)_K=(\T'_\m)_K\times(\T''_\m)_K$ (new and old
  eigensystems are disjoint, as the $U_q$-eigenvalues on the $q$-old part have absolute value $\sqrt q$ and those on the $q$-new part
  absolute value $1$);
\item $\T'_\m\cong\T_{B,\m}$ via $\theta_s$ (\pathtt{if:s01-targets}), and $e\,\Zp[S_s]_\m$ is free of rank one over $\T_{B,\m}$ with
  $\T_{B,\m}$ Gorenstein (16.8, [Hs, Thm.~4.5]).
\end{enumerate}
Then both sides of $\HypXiso(s)$ are free of rank one over $\T'_\m$, hence isomorphic (non-canonically). \stpmc.\\\emph{Lean:} the algebraic half (for $H$ free of rank one over a local $A$-Gorenstein $T$, with $T/\fa$ Gorenstein and the hypotheses of Prop.~18.25, $H[\fa]\cong T/\fa$) is \stformal{} (\pathtt{Comparison/XIso.lean}, sorry-free): \lean{Langlandscape.Comparison.xiso\_of\_gorenstein}; the verification of (i)--(iii) for the Hecke algebras of the modular curve and the $\Lam$-adic compatibility are not formalized.
Compatibility in $s$ (choosing compatible generators along the tower) additionally needs $\Lam$-adic freeness of $\varprojlim_sH$ over
the $\Lam$-adic $\T_\m$ (Hida, [Hi86a], not in \pathtt{literature/}) and of ${\Mord}_\m$ (16.8); \stsketch.
\emph{Freeness of rank one} is Hsieh's Thm.~4.5, which also assumes his hypothesis (CR) and the support condition on $N^-$ (to be carried explicitly in (i), (iii)); Helm's Cor.~8.11 and Rem.~8.12 are multiplicity-one (dimension) statements, entering via Hsieh's proof.\\
References: \lit{hsieh\_2017\_hida-families-and-p-adic-triple-product-l-functions\_arxiv}{Thm.~4.5, p.~34, whose proof cites
[Wil95, Thm.~2.1] and [Hel07, Cor.~8.11, Rem.~8.12]};
\lit{helm\_2004\_maps-between-modular-jacobians-and-jacobians-of-shimura-curves\_arxiv}{Cor.~8.11, Remark~8.12, p.~31}.
\end{quoted}

\begin{quoted}{Proposition}{18.25}{151}{prop:s05-subquot}{$q$-new sub- versus quotient lattice; seed Prop.~subquot; \pathtt{if:s05-subquot}}{chapters/s05/subquot.tex}
Let $A$ be a Noetherian local domain with fraction field $K$, $\T$ a local, finite free $A$-algebra, and $\fa,\fb\subseteq\T$ ideals with
$\T/\fa$, $\T/\fb$ $A$-torsion-free, $\fa\cap\fb=0$ and $\T_K=(\T/\fa)_K\times(\T/\fb)_K$ (equivalently, by \pathtt{NewOldData.split\_iff}:
there is $c\in A\setminus\{0\}$ with $c\in\fa+\fb$). Put $\T':=\T/\fa$, $\T'':=\T/\fb$ and $\cc_{\T'}:=(\fa+\fb)/\fa\subseteq\T'$. Let $H$ be free
of rank $r\ge0$ over $\T$, and $H[\fa]:=\{h\in H:\fa h=0\}$. Then
\[H[\fa]=\fb H\cong\fb^{\oplus r},\qquad H/\fa H\cong\T'^{\oplus r},\]
the natural map $H[\fa]\hookrightarrow H\twoheadrightarrow H/\fa H$ is injective, and its cokernel is $(\T'/\cc_{\T'})^{\oplus r}\cong(\T/(\fa+\fb))^{\oplus r}$.
In particular the cokernel is killed by $c$. \stproved{} (below), using only Lemma~D(a) ($\Ann_\T\fa=\fb$; \texttt{if:s06-lemD}; there
\stproved).
\emph{Lean:} \stformal{} (\pathtt{CommAlg/SubQuot.lean}, sorry-free, standard axioms only): \lean{Langlandscape.CommAlg.subquot\_kernel}, \lean{Langlandscape.CommAlg.subquot\_coker}.
\end{quoted}

\begin{quoted}{Proposition}{18.20}{144}{prop:s05-Xlambda}{$\Lam$-adic version; seed Prop.~Xlambda}{chapters/s05/hypX.tex}
Assume $\HypXchar(M;q)$. All limits are over $s\ge1$ (at $s=0$ there is no $U_p$ in $\cH_0$, hence no $e$; $s=0$ plays no role).
\begin{enumerate}[label=(\alph*)]
\item The maps $\beta_{s's}\otimes\Zp$ restrict to $e(\Xq(J_{s'})\otimes\Zp)\to e(\Xq(J_s)\otimes\Zp)$, and the $\varphi_s$ induce an isomorphism
  \[\Phi:\varprojlim_{s\ge1}e\big(\Xq(J_s)\otimes\Zp\big)\xrightarrow{\sim}\Mord(M),\]
  compatible with $T_\ell\mapsto T_\ell$ ($\ell\nmid Mpq$), $U_\ell\mapsto U_\ell$ ($\ell\mid M$), $U_p\mapsto U_p$, $U_q\mapsto\Pi_q$, the classical
  $\langle d\rangle\mapsto$ the adelic $\langle d\rangle^{-1}$ for $d\in(\Z/M)^\times\times\mu_{p-1}$ (embedded in $(\Z/Mp^s)^\times$ via $d\bmod M$ and the
  Teichmüller lift mod $p^s$; the finite group of \pathtt{def:s04-Mord}), and $\iota$-semilinear for $\Lam$: $\Phi(\lambda x)=\iota(\lambda)\Phi(x)$.
  In particular $\iota$ carries the seed's weight-$k$ point $P_k=\ker([u]\mapsto u^{k-2})$ to $P^{\rm ad}_k=\ker([u]\mapsto u^{2-k})$ of stage 04.
  \stproved{} given 13.4(a).
\item The left side of (a) is free over $\Lam$ of rank $r=\rk_{\Zp}e\,\Zp[S_1]$, and for $s\ge1$ its $\iota$-twisted specialisation
  $(\varprojlim e(\Xq(J_\bullet)\otimes\Zp))\otimes_{\Lam,\iota}\Zp[\Delta_s]$ is $e(\Xq(J_s)\otimes\Zp)$, Hecke-equivariantly
  ($\Delta_s$ the image of $1+p\Zp$ in $(\Z/p^s)^\times$). \stpmc{} (transport of Thm.~control(a) of stage 04, 16.8).
\item (Twisted pairings.) For $s\ge1$ let $E_s:=e(\Xq(J_s)\otimes\Zp)$ ($U_p$ is invertible on $E_s$), $w=w_{Mp^s}$, and
  \[[x,y]_s:=u_s(x,\,w\,U_p^{-s}y)\qquad(x,y\in E_s).\]
  \begin{enumerate}[label=(c\arabic*)]
  \item $[\,,\,]_s$ is perfect and $\cH_s$-equivariant: $[hx,y]_s=[x,hy]_s$ for $h\in\cH_s$. \stpmc{} (batch 6: from Lemma~18.19, (R2), (X1), (X2) and
    13.4(d)).
  \item \emph{(Expected; not proved.)} Through $\varphi_s$, $[\,,\,]_s$ corresponds to the twisted Brandt pairing of stage 04 (Hsieh's $B_{N,s}$,
    \texttt{def:s04-BN}), compatibly along the tower. \stopen{} (expected: a heuristic with no precise statement yet; the label \texttt{conjectural} is used only for want of a weaker label). This is not yet a precise statement: the normalisation of the comparison
    (which modules, which constant) is item (vi) of Remark~18.22, which lists the missing steps (i)--(vi). Nothing in this chapter
    uses (c2) as proved.
    \provisional{S05-M10: (c2) downgraded from \stsketch{} to \stconj/\stopen, after S05-R1 (batch-6 re-read, GAP: the norm twist was missing from
    the list). The list is completed in Remark~18.22; Lemma~18.21 locates the norm twist on the Brandt side.
    Downstream uses flagged: the $\Hom_\Lam(\Mord,\Lam)(1)$-form of (d2), the sentence after the proof, Remark~18.36, and the
    $\eta^B_F$ row of the table in \S18.4.1.}
  \end{enumerate}
  \provisional{batch 6 (E3): the seed and stage tex twist by $U_p^{-s}w_{p^s}$. Two changes, both corrections.
  (i) We use $w_{Mp^s}$, not $w_{p^s}$. Stage 04's $\tau^D_{Np^\alpha}$ has components at all primes of $M$ (\texttt{def:s04-BN}, translation (ii)).
  Also, for the partial $w_{p^s}$, Atkin--Li Prop.~1.3 gives $w_{p^s}U_\ell w_{p^s}^{-1}=\langle\ell\rangle_{p^s}^{-1}U_\ell$ for $\ell\mid M$. That is not
  $U_\ell^\Ros$ on the $\ell$-old part, where $U_\ell$ is not normal. So the (R1)-analogue that (c1) needs fails for $w_{p^s}$ when $M>1$ (\stsketch).
  (ii) We apply $U_p^{-s}$ first, then $w$. With the other order, equivariance for $h=U_p$ would need $U_p^\Ros$ to commute with $U_p$, which fails in general.
  The order agrees with $B_{N,\alpha}$, where $U_p^{-\alpha}$ is applied before translating by $\tau^D$.}
\item (Monodromy filtration of the ordinary tower.) Assume also $\HypXgal$. Let
  \[\Tord:=\varprojlim_{s\ge1}e_*\,T_pJ_s,\]
  the limit along the Albanese pushforwards $\pi_*:J_{s'}\to J_s$. Here $e_*$ is the ordinary projector of the Albanese $U_{p,*}=U_p^\Ros$, which commutes
  with $\pi_*$ (Lemma~18.9(c)). Each $h\in\cH_s$ acts through its Albanese version $h_*:=h^\Ros$, and $\Lam$ through the Albanese diamonds.
  This is Hida's convention
  \lit{hida\_1986\_galois-representations-into-gl2-zp-x-attached-to-ordinary-cusp-forms}{\S8, pp.~595--596: $J_r$ viewed as
  $\mathrm{Alb}(X_r)$, covariant Hecke action, limit along the covariant maps, $e$ attached to $T(p)$}.
  The filtrations $T_t\subset T^{I_q}\subset T_pJ_s$ (Thm.~18.30(v)) induce a $\Lam[G_{\Q_q}]$-stable filtration
  $\Tord_t\subset{\Tord}^{I}\subset\Tord$ with:
  \begin{enumerate}[label=(d\arabic*)]
  \item \emph{top:} $\Tord/{\Tord}^{I}\cong\varprojlim_se(\Xq(J_s)\otimes\Zp)$, the left side of (a); hence $\Tord/{\Tord}^I\cong\Mord(M)$ via $\Phi$.
    Under this, $h_*\mapsto\theta(h)$ for the operators $h$ listed in (a), and the $\Lam$-structure is $\iota$-semilinear as in (a). $\Frob_q$ acts as $\Pi_q^{-1}$
    (the Frobenius clauses of (d1), (d2) come from the Frobenius clause of (X4); FD1 is decided (Fabian, 2026-11-15): the model of $X_s$ is the $\mu_N$-model with Picard functoriality.).
  \item \emph{bottom:} $\Tord_t\cong\Mord(M)\otimes\Zp(1)$, with the same Hecke and $\Lam$-structure. $\Frob_q$ acts as $q\,\Pi_q^{-1}$.
  \item \emph{monodromy:} under (d1), (d2), the monodromy operator $\Tord/{\Tord}^I\to\Tord_t(-1)$ is $\pm$ the identity of $\Mord(M)$, with $+$ for the
    identifications fixed in the proof.
  \item \emph{middle:} ${\Tord}^I/\Tord_t\cong\Tord(M)^{\oplus2}$, Hida's ordinary Tate module at level $M$, via the two degeneracy maps.
    It is $\Lam$-free \lit{hida\_1986\_galois-representations-into-gl2-zp-x-attached-to-ordinary-cusp-forms}{Thm.~3.1(ii), p.~562
    ($V^{\rm ord}$ is $\Lam$-free of finite rank), with $J^{\rm ord}_\infty\cong V^{\rm ord}$, p.~596}. Hence $\Tord$ is $\Lam$-free.
  \end{enumerate}
  Status: (d1)--(d3) \stpmc{} (batch 6). The Hecke and transition structure of (d1), (d2) uses only $\HypXchar$; the Frobenius and $N$ use the
  corresponding clauses of (X4). (d4) \stsketch{} (missing step at the end of the proof).
  \provisional{batch 6 (E3): statement of (d) changed. Batch 5 found that the projector and the tower were not fixed, and that the pieces had been computed
  with the wrong (Picard) projector. Now:
  \begin{itemize}
  \item $\Tord$ is Hida's Albanese tower with $e_*$.
  \item The bottom piece is $\Mord(M)(1)$, identified canonically through the perfect $(e,e^*)$ Brandt pairing. It replaces the former
    $\Hom_\Lam(\Mord,\Lam)(1)$; identifying the two needs the $\Lam$-adic twisted pairing of (c2), which is \stopen{} (\stconj; Remark~18.22),
    so the $\Hom_\Lam$-form is not claimed.
  \item The top piece is $\Mord(M)$ with $h_*\mapsto\theta(h)$. This is the former ${\Mord}^\Ros$ read through $h_*=h^\Ros$: the same Rosati twist as in (X4).
  \item Hida's $\Lam$-freeness is now needed only for the middle piece, and its page is checked.
  \end{itemize}
  The status is raised from \stsketch{} to \stpmc{} for (d1)--(d3).}
\end{enumerate}
\par\smallskip\noindent{\scriptsize\itshape [Status paragraph following the statement in the chapter:]}\par
\noindent\emph{Lean status} (part (a), finite level only): \stformal{} \emph{conditional on} $\HypXchar$ (carried as the argument \texttt{hX : Hyp.X.char}), for an abstract \texttt{CharGroupData} $G$ and not for $\Xq(J_s)$ itself (\lean{Langlandscape.Comparison.xlambda\_finite}, \pathtt{Comparison/XLambda.lean}): for each $s\ge1$ it gives $\mathtt{Nonempty}$ of a $\Zp$-linear equivalence $e(\Zp\otimes G.X_s)\simeq e\,\Zp[S_s]$, with $e$ the ordinary part of $U_p\otimes1$ resp.\ of $U_p$. It is a corollary of \lean{Langlandscape.Comparison.xlambda\_finite\_iso} (under $\HypXchar$, a $\Zp$-linear isomorphism $\Psi:\Zp\otimes G.X_s\simeq\Zp[S_s]^0$ intertwining $U_p\otimes1$ and $U_p$). The Lean statement is only an existence (rank-level) statement: that the equivalence is \emph{induced by} $\varphi_s$ is not part of the \texttt{Nonempty} statement (it is for the definition \texttt{xlambdaOrd}; the open point R2-2 of the Lean audit). The inverse limit over $s$, parts (b)--(d) and the Galois and monodromy statements are not formalized (\leanplanned{}).
\end{quoted}

\begin{quoted}{Conjecture}{18.40}{159}{conj:s05-O2}{O2, a simplicial realisation; open}{chapters/s05/simplicial.tex}
There is a Hecke-equivariant map of finite simplicial (orbi-)sets, defined over $\Z[\frac1{6}]$ at most, from $\GDR_q(U_s)$ to a simplicial
model of the special fibre of $X_s$ at $q$ (or to a $q$-adic tubular neighbourhood), inducing $\varphi_s^{-1}$ on $H_1$. \stconj{};
equivalently (by Remark~18.38) a degree-shifting map between the two $S$-arithmetic complexes. The most promising concrete
form is the Deuring map on supersingular points, which exists only after reduction mod $q$; we expect no characteristic-zero Betti
version. \stopen.
\end{quoted}
